\uuid{urVv}
\titre{Résolution d'équation différentielle du premier ordre}
\auteur{bodin ; croizat ; sacre}
\organisation{exo7}
\licence{CC BY-NC-SA 4.0}
\datecreate{2022-03-08}
\niveau{L1}
\module{Analyse}
\chapitre{Equation différentielle}
\sousChapitre{Résolution d'équation différentielle du premier ordre}
\difficulte{1}
\theme{}
\exo7qcmid{hilisit.equations-differentielles.038}

\contenu{
\begin{question}
\qclassification{225.01}
Soit l'équation différentielle $2y'+4y=3$. Quelles sont les affirmations vraies ?
\begin{answers}
  \bad{La seule solution constante est $y=3/2$.}
  \good{La seule solution constante est $y=3/4$.}
  \bad{Les solutions sont $C e^{-4x}-3$  avec $C$ constante réelle.}
  \good{Les solutions sont $C e^{-2x}+3/4$  avec $C$ constante réelle.}
\end{answers}
\begin{explanations}
La seule solution constante est $y=3/4$ : c'est ce qu'on retrouve dans l'équation différentielle lorsqu'on cherche $y$ constante avec donc $y'=0$ : l'équation devient $2y = 3/2$ donc $y = 3/4$.
On peut réécrire l'équation différentielle $y'=-2y+3/2$, dont les solutions sont $C e^{-2x}+3/4$ avec $C$ constante réelle.
\end{explanations}
\end{question}
}
